\centerline{\bf \Huge 10 класс}

\problem{1}
Каждую грань кубика разбили на четыре одинаковых квадрата, а затем раскрасили эти квадраты в несколько цветов так, что квадраты, имеющие общую сторону, оказались окрашенными в различные цвета. Какое наибольшее количество квадратов одного цвета могло получиться?

\problem{2}
Внутри параллелограмма $A B C D$ выбрана точка $E$ так, что $A E=D E$ и $\angle A B E=90^{\circ}$. Точка $M$ - середина отрезка $B C$. Найдите угол $D M E$.

\problem{3}
Множество состоит из целых чисел, причем в нем есть и положительные, и отрицательные числа. Если числа $a$ и $b$ содержатся в множестве, то в нём также содержатся числа $2 a$ и $a+b$. Докажите, что это множество содержит разность любых двух своих элементов.

\problem{4}
$a, b, c$-стороны треугольника с периметром 1. Докажите неравенство

$$
\frac{1+a}{1-2 a}+\frac{1+b}{1-2 b}+\frac{1+c}{1-2 c} \geqslant 12 .
$$
 
\problem{5}
Дана таблица $100 \times 100$. Для любого $k(1 \leqslant k \leqslant 100)$, в $k$-й строке таблицы записаны числа $1,2, \ldots, k$ в возрастающем порядке слева направо, но не обязательно в последовательных клетках; остальные $100-k$ клеток заполнены нулями. Докажите, что найдутся два столбца, сумма чисел в одном из которых хотя бы в 19 раз превосходит сумму чисел в другом.